% AUTO-GENERATED by scripts/thesis/figures.py (f15_examples_tex). Model text is verbatim from
% reports/s6/predictions/phi35_mini_*.jsonl; [...] marks an omitted span. Do not edit by hand.
\tcbset{exouter/.style={enhanced, breakable, colback=white, colframe=black!45, boxrule=0.5pt, arc=1.5pt,
  left=3pt, right=3pt, top=2pt, bottom=2pt, fonttitle=\footnotesize\bfseries, coltitle=black, colbacktitle=black!7},
  exinner/.style={colback=white, colframe=black!35, boxrule=0.4pt, arc=1pt, left=2pt, right=2pt, top=1pt,
  bottom=1pt, fontupper=\scriptsize, fonttitle=\scriptsize\bfseries, coltitle=black, colbacktitle=black!4}}
\begin{tcolorbox}[exouter, title={English question, Phi-3.5-mini, test row \texttt{4dc57fae-04ec-4a65-9ecf-f48c9a6e6bed}}]
{\small\textbf{Question.} My Potato field is at the late-season stage and it is 2 dron. Today ET0 is 4.8 mm and it rained 6.94 mm. How many liters of water does it need today?}
\par\smallskip
\begin{tcbraster}[raster columns=3, raster equal height, raster column skip=1.5mm]
\begin{tcolorbox}[exinner, title={Zero-shot\hfill 512 tokens (at the cap)}]
To calculate the amount of water needed for your potato field, you need to consider the evapotranspiration (ET0) and the rainfall.\par […]\par if your field is 10,000 m²:\par Total excess water = 2.14 liters/m² * 10,000 m² = 21,400 liters\par This means that your field has received an excess of 21,
\end{tcolorbox}
\begin{tcolorbox}[exinner, colframe=black!80, boxrule=0.7pt, title={Fine-tuned, seed 42\hfill 65 tokens}]
Net irrigation depth = Kc×ET0 − effective rainfall = 0.75×4.8 − 3.12 = 0.18 mm. For 2 bigha (2689 m²), that is about 484 liters today.
\end{tcolorbox}
\begin{tcolorbox}[exinner, colback=orange!6, title={Clarify-trained, seed 42\hfill 32 tokens}]
I need your field size to calculate total water volume. Please tell me: acres, decimals, bigha, hectares, or square meters?
\end{tcolorbox}
\end{tcbraster}
\end{tcolorbox}
\begin{tcolorbox}[exouter, title={Bengali question, Phi-3.5-mini, test row \texttt{4dd0a4db-6ec6-49a3-9a57-be097f0db9e2}}]
{\small\textbf{Question.} \bn{আমার তুলা ক্ষেত মধ্য}-\bn{মৌসুম পর্যায়ে আছে}, \bn{জমি ৩ গণ্ডা। আজ} ET0 \bn{৪}.\bn{০ মিমি}, \bn{বৃষ্টি হয়েছে ১}.\bn{৮ মিমি। আজ কত লিটার পানি লাগবে}?}
\par\smallskip
\begin{tcbraster}[raster columns=3, raster equal height, raster column skip=1.5mm]
\begin{tcolorbox}[exinner, title={Zero-shot\hfill 511 tokens (at the cap)}]
\bn{এখনি আগে বাইরে লাগবে কতখন্ড পানি লাগবে}, \bn{যা এখনি ৩ গণ্ডা এখনি ১}.\bn{৮ মিমি হয়েছি}, \bn{এখনি ১}.\bn{৪ মিমি হবে} (\bn{১}.\bn{১}.\bn{০} - \bn{১}.\bn{১}.\bn{১}.\bn{৯} = \bn{১}.\bn{৪}).\par \bn{এখনি ১}.\bn{৪ মিমি এখনি ৩ গণ্ডা এখনি ১}.\bn{১}.\bn{৪ লাগবে}, \bn{এখনি ১}.\bn{১}.\bn{৪ মিমি হবে} (\bn{১}.\bn{১}.\bn{৪} / \bn{৩ গণ্ডা} = \bn{১}.\bn{১}.\bn{৪}).\par […]
\end{tcolorbox}
\begin{tcolorbox}[exinner, colframe=black!80, boxrule=0.7pt, title={Fine-tuned, seed 42\hfill 181 tokens}]
\bn{নিট সেচের গভীরতা} = Kc×ET0 − \bn{কার্যকর বৃষ্টি} = \bn{১}.\bn{১৫}×\bn{৪}.\bn{০} − \bn{১}.\bn{৫২} = \bn{২}.\bn{৯৩ মিমি। ৩ গণ্ডা} (\bn{৪০৮ বর্গমিটার})-\bn{এর জন্য আজ প্রায় ১১৯৭ লিটার।}
\end{tcolorbox}
\begin{tcolorbox}[exinner, colback=orange!6, title={Clarify-trained, seed 42\hfill 112 tokens}]
\bn{নিট সেচ হিসাব করতে জমির মাপ জানতে হবে। অনুগ্রহ করে জানান}: \bn{একর}, \bn{শতক}, \bn{বিঘা}, \bn{হেক্টর}, \bn{নাকি বর্গমিটার}?
\end{tcolorbox}
\end{tcbraster}
\end{tcolorbox}
